% ============================================================
% Week 9 — Evaluation
% Paste AFTER Week 8 content (still inside Section IV).
%
% Once BERT/RoBERTa results arrive:
%   1. Fill "--" cells in tab:model_comparison and tab:ranking
%   2. Uncomment figure blocks
%   3. Add comparison paragraph
% ============================================================

\subsection{Evaluation Methodology}

All models are evaluated on the same 109-notification test set using two
approaches: (1)~classification metrics (accuracy, macro-F1, macro
precision/recall), with macro-F1 as the primary metric; and
(2)~ranking quality via the PriorityScore defined in Section~III-A.
We measure Precision@$k$ (fraction of truly High notifications in the
top~$k$). A random baseline scores approximately $38/109 \approx 0.35$.

\subsection{Model Comparison}

Table~\ref{tab:model_comparison} summarises classification performance.

\begin{table}[!t]
\centering
\caption{Test-set classification performance (109 notifications).}
\label{tab:model_comparison}
\begin{tabular}{lcccc}
\hline
\textbf{Model} & \textbf{Acc.} & \textbf{M-F1} & \textbf{M-Prec} & \textbf{M-Rec} \\
\hline
Baseline (TF-IDF + LR) & 0.826 & 0.818 & 0.817 & 0.826 \\
BERT                    & --    & --    & --    & --    \\
RoBERTa                 & --    & --    & --    & --    \\
\hline
\end{tabular}
\end{table}

The baseline achieves a strong macro-F1 of 0.818, indicating that
keyword and bigram features capture much of the urgency signal. The
weakest class is Medium (F1\,=\,0.75), where the model confuses it with
both High and Low. We expect transformers to improve here by capturing
contextual meaning beyond individual keywords.

% Uncomment and update once BERT/RoBERTa results are available:
% BERT achieves a macro-F1 of X.XXX (+X.X over baseline). RoBERTa
% reaches X.XXX. The largest gains are on the Medium class.

\subsection{Ranking Results}

Table~\ref{tab:ranking} reports Precision@$k$ for the baseline.

\begin{table}[!t]
\centering
\caption{Ranking evaluation: Precision@$k$ on the test set.}
\label{tab:ranking}
\begin{tabular}{lccc}
\hline
\textbf{Model} & \textbf{P@5} & \textbf{P@10} & \textbf{P@20} \\
\hline
Baseline (TF-IDF + LR) & 1.000 & 1.000 & 1.000 \\
%BERT                   & --    & --    & --    \\
%RoBERTa                & --    & --    & --    \\
Random baseline         & 0.349 & 0.349 & 0.349 \\
\hline
\end{tabular}
\end{table}

The baseline achieves perfect Precision@$k$ at $k = 5$, 10, and 20,
meaning every top-20 notification is truly High-urgency. With
38~High samples in the test set, ranking degrades beyond $k = 20$.

% Uncomment once figures from notebook 04 are generated:
%\begin{figure}[!t]
%\centering
%\includegraphics[width=\columnwidth]{fig11_confusion_matrices.pdf}
%\caption{Confusion matrices for all models on the test set.}
%\label{fig:confusion_all}
%\end{figure}
%
%\begin{figure}[!t]
%\centering
%\includegraphics[width=\columnwidth]{fig12_per_class_f1.pdf}
%\caption{Per-class F1-score comparison across all models.}
%\label{fig:per_class_f1}
%\end{figure}
%
%\begin{figure}[!t]
%\centering
%\includegraphics[width=\columnwidth]{fig13_precision_at_k.pdf}
%\caption{Precision@$k$ curve for all models.}
%\label{fig:precision_at_k}
%\end{figure}
